\chapter{Att kontrollera ett påstående}
\label{app:verifiera}

Det här materialet gör, som allt undervisningsmaterial, en rad påståenden.
De flesta går att pröva vid tangentbordet: skriv av programmet, kör det och
se efter.  Men några handlar om vad forskningen säger --- om hur nybörjare
faktiskt tänker om slingor, och om hur man bäst lär ut dem --- och sådana
påståenden går inte att köra.  De måste kontrolleras mot litteraturen.

Samma fråga möter dig varje gång du läser en lärobok, en webbsida eller ett
svar från en språkmodell, och varje gång du själv ska skriva en rapport.
Den här bilagan beskriver en arbetsgång för att kontrollera ett påstående,
och \cref{tab:paastaenden} visar var varje påstående i det här materialet
kontrolleras.  Bilagan som följer är arbetsgången genomförd, och går att
läsa fristående från resten.

\begin{table}[htbp]
  \centering
  \caption{Materialets påståenden om forskning, och var de kontrolleras.}
  \label{tab:paastaenden}
  \begin{tabular}{@{}p{0.52\linewidth}p{0.42\linewidth}@{}}
    \toprule
    Påstående & Kontrolleras \\
    \midrule
    Att det är bra att inte upprepa sig i ett program --- DRY --- och att
      principen kommer från Hunt och Thomas
      & Redan belagt i \emph{Algoritmiskt tänkande}, bilaga E \\
    Att en funktion bör ha ett enda ansvar --- SRP
      & Redan belagt i \emph{Funktioner}, bilaga B \\
    Att kontrast måste upplevas samtidigt för att en aspekt ska urskiljas
      & Andrahandskälla, se kommentaren efter tabellen \\
    Att nybörjare håller de missuppfattningar om slingor som materialet
      undervisar mot & \Cref{app:missuppfattningar} \\
    Att det lönar sig att låta studenten försöka själv innan förklaringen ges
      & Redan belagt i \emph{Övning: Funktioner och variabler}, bilaga~B \\
    \bottomrule
  \end{tabular}
\end{table}

Några rader förtjänar en kommentar.  \emph{SRP} namnges i
\cref{sec:inmatning} men prövas inte om här --- kontrollen hör hemma i
föreläsningen \emph{Funktioner}, där principen införs, och den är
genomförd där, i bilaga B.  Sökningen om att försöka först gjordes för
det här materialet, men kapitlet som redovisar den står i \emph{Övning:
Funktioner och variabler}, bilaga~B, eftersom övningen är det första
material i kursen som gör påståendet.  Raden om kontrast gäller
påståendena om lärande i marginalens lärarnoter, som vilar på Martons
\emph{Necessary Conditions of Learning}; citaten därifrån är
kontrollerade mot en andrahandsuppteckning och inte mot boken, vilket är
antecknat i litteraturlistans källfil.  Poängen med att de \emph{står} i
tabellen är enkel: det som inte är kontrollerat här ska räknas upp lika
noga som det som är det.

\section{Gör påståendet till en fråga}

Ett påstående går att kontrollera först när det är tydligt vad som skulle
kunna visa att det är \emph{fel}.  Formulera det därför som en fråga med ett
svar: \enquote{tror nybörjare att kod intill en slinga hör till den?} kan
besvaras med ja eller nej, medan \enquote{slingor är svåra} inte kan det.
Ingen tänkbar källa skulle visa att det senare är fel, eftersom
\enquote{svårt} är ett omdöme och inte ett faktum.  Den som säger så menar
oftast något annat, till exempel \enquote{gör nybörjare fler fel på slingor
än på förgreningar?} --- och det är en fråga som kan få svaret nej, och
därför kan prövas.

Ofta döljer sig flera påståenden i ett.  \enquote{Principen DRY är bra och
kommer från Hunt och Thomas} är två frågor --- \emph{håller den?} och
\emph{varifrån kommer den?} --- och de kräver olika slags källor.

\section{Välj rätt sorts källa}

Olika frågor besvaras av olika källor.
\begin{description}
  \item[Uppslagsverk] för vad ett ord \emph{betyder}.
  \item[Primärkällan] för var en idé \emph{kommer ifrån}.  Vill man veta om
    DRY härrör från Hunt och Thomas läser man deras bok, inte en lärobok som
    återger den.  Finns flera utgåvor väljer man originalet.
  \item[Översikter] (\foreignlanguage{english}{reviews}) för om något är
    \emph{etablerat}, eller för vad forskningen \emph{sammantaget} säger.  En
    systematisk översikt går igenom hundratals studier efter en angiven metod
    och väger därför tyngre än en enskild studie.
  \item[Enskilda studier] för ett specifikt resultat --- med förbehållet att
    en enskild studie bara visar att något observerats en gång.
\end{description}
Var man hittar dem: universitetsbibliotekets söktjänst, och databaser som
Scopus, Web of Science, IEEE Xplore, ACM Digital Library, DBLP (datalogi) och
OpenAlex (OA; öppen, och inte samma sak som
\foreignlanguage{english}{open access}).  Google Scholar täcker brett men
rankar ogenomskinligt.  I bilagorna som följer användes
kommandoradsverktyget \texttt{scholar}\footnote{Paketet \texttt{scholarcli} på
  PyPI: \url{https://pypi.org/project/scholarcli/}.}, som ställer samma fråga
till flera databaser på en gång och sparar allt som hände.  Det du får ut av
bibliotekets söktjänst är detsamma, bara långsammare.

\section{Sök så att du kan ha fel}

Den som söker på namn hon redan känner till hittar det hon väntar sig.  Sök
därför \emph{ämnesdrivet} --- på begrepp, inte författare.  De källor som
bara går att nå som kända dokument, som äldre rapporter utan
\textsc{doi}, anges som sådana, med skälet.

Några praktiska regler:
\begin{itemize}
  \item Håll frågorna korta.  Flera korta frågor slår en lång.
  \item Ställ \emph{samma} frågor i alla databaser, och sök begreppets alla
    namn.  Varje databas har dessutom sin egen syntax: Web of Science vill ha
    \texttt{TS=(\dots)}, Scopus \texttt{TITLE-ABS-KEY(\dots)}, medan DBLP bara
    tar en ordlista där alla ord måste matcha.
  \item Sök också efter \emph{invändningar}: lägg till ord som
    \enquote{critique}, \enquote{limitations} eller \enquote{failure}.  Ett
    påstående som överlevt en motsökning är värt mer än ett som aldrig
    prövats, och en invändning som hittas blir ett \emph{förbehåll} i svaret,
    inte ett nederlag.
  \item Räkna med att databasernas rankning missar somligt.  Då hämtas källan
    som \emph{känt dokument}: man vet att den finns och slår upp den direkt.
    Det är i sin ordning, bara det sägs.
\end{itemize}

\section{Läs i fulltext, och skriv ner hur}

En träff i en databas är inte ett belägg.  Sammanfattningen
(\foreignlanguage{english}{abstract}) räcker inte heller: den säger vad
författarna vill ha sagt, inte vad de visat.  Läs hela texten och leta upp
\emph{det ställe} som styrker påståendet --- ett ordagrant citat med
sidnummer.  Hittar du inget sådant ställe stöder källan inte påståendet, hur
passande titeln än är.  Räkna aldrig en källa du inte läst.

Det gäller också en källa du ärvt.  Att ett påstående står med en fotnot i ett
annat material betyder inte att du har kontrollerat det; du har kontrollerat
att någon annan påstod det.  Bilagorna som följer bygger delvis på källor från
kursens forskningsunderlag om missuppfattningar
\autocite{SooriBosk2026}, och varje sådan källa har
lästs om här --- vilket i två fall ändrade vad materialet får säga
(\cref{sec:miss-diskussion}).

Skriv sedan ner hur du gick till väga, så att någon annan kan följa spåret.
I det här materialet står den anteckningen som en kommentar ovanför varje post
i litteraturlistans källfil, med fasta fält:
\begin{description}
  \item[\texttt{CLAIM}] vilket påstående källan ska styrka;
  \item[\texttt{FOUND-VIA}] hur den hittades --- sökfråga och databas, eller
    \enquote{känt dokument};
  \item[\texttt{PICKED}] varför just den, framför andra träffar;
  \item[\texttt{QUOTE}] det ordagranna citatet, med sida;
  \item[\texttt{VERIFIED}] att fulltexten lästs, och var;
  \item[\texttt{COUNTER}] vad motsökningen gav;
  \item[\texttt{DATE}] när.
\end{description}
Det ser byråkratiskt ut och tar ett par minuter per källa.  Det är också det
enda som gör det möjligt att om ett år svara på frågan \enquote{hur vet du
det?}

\section{Dra slutsatsen --- med förbehåll}

Svaret skrivs ut tillsammans med det som talar emot.  Förbehållen är inte
svagheter i svaret; de \emph{är} svaret, och de avgör hur påståendet får
användas i materialet.  Ett påstående vars belägg bara räcker till
sammanfattningsnivå ska formuleras svagare än ett som är läst i fulltext, och
det ska synas att så är fallet.  Varje kapitel slutar därför med ett
avsnitt \emph{Fortsatt arbete}: vad just den sökningen lämnade ogjort,
och vad som skulle behöva undersökas för ett säkrare svar.

\section{En anmärkning om språkmodeller}

Den här föreläsningens egen bilaga gjorde ingen ny sökning
(\cref{sec:miss-metod}), men sökningen om att försöka först --- redovisad i
\emph{Övning: Funktioner och variabler}, bilaga~B --- gav hundratals
träffar.  En språkmodell användes för att \emph{sortera} dem --- hör
träffen till ämnet över huvud taget, och åt vilket håll pekar den? --- så att gallringen kunde granskas.
Lägg märke till arbetsfördelningen: modellen sorterade, men varje källa som
citeras ska läsas och väljas av en människa, och även sorteringen ska
bekräftas av en människa.  Modellens klassning står därför med sin konfidens
i tabellerna, så att den kan ifrågasättas.  Använd gärna språkmodeller för
att hitta och sortera; låt dem aldrig vara det sista ledet.  De kan göra fel,
men det är du själv som måste stå till svars för innehållet.

\chapter{Håller missuppfattningarna om slingor?}
\label{app:missuppfattningar}
\chapterprecis{Författaren har ännu inte granskat resultaten i den här
  bilagan i sin helhet.}

Materialet undervisar mot sex konkreta missuppfattningar: att det som
upprepas är den text som syns i terminalen snarare än slingans kropp
(\cref{sec:for}); att kod \emph{intill} en slinga hör till den
(\cref{sec:indrag}); att vad som helst kan stå efter
\mintinline{python}{in} (\cref{sec:generalisering}); att en räknare inte
behöver sättas innan slingan börjar (\cref{sec:while}); att en
\mintinline{python}{while}-slinga tar slut även om ingenting i kroppen
ändrar villkoret (\cref{sec:while}); och att en upprepning gör samma nytta
som en förgrening när handlingen bara ska utföras en gång
(\cref{sec:upprepa-eller-valja}).  Frågan är: \emph{är det
verkligen missuppfattningar som nybörjare håller, eller är de bara rimliga
att tänka sig?}

\section{Metod}
\label{sec:miss-metod}

Ingen ny sökning gjordes.  Påståendena täcks av kursens eget
forskningsunderlag om missuppfattningar i programmering
\autocite{SooriBosk2026}, vars systematiska
sökning ställde bland annat frågan \enquote{loop iteration misconceptions
novice programming} till sju databaser --- OpenAlex, DBLP, arXiv, Web of
Science, Scopus, IEEE Xplore och Semantic Scholar --- och sedan gallrade
träffarna med en språkmodell och för hand.  Att göra om den sökningen här
skulle inte tillföra något; däremot duger det inte att bara ärva fotnoterna.
Därför lästes varje källa som citeras här om i fulltext, och varje citat
söktes upp på nytt i källan (\cref{app:verifiera}).  Det gav både
bekräftelser och två ändringar i materialet: resultatet står i
\cref{sec:miss-resultat} och ändringarna i \cref{sec:miss-diskussion}.

\section{Resultat}
\label{sec:miss-resultat}

\paragraph{Kod intill slingan.}  Den bäst belagda av de sex.
\textcite{Sleeman1984} intervjuade gymnasieelever som skrev och felsökte
Pascalprogram, och noterade felet att en \mintinline{text}|WRITELN| skriven
bredvid slingan lästes som en del av den:
\textcquote[15]{Sleeman1984}{WRITELN adjacent to the loop is included in it
\ldots{} This error occurred frequently and was noted with nearly half the
students interviewed}.\footnote{Rapporten finns bara som en inskannad
  text från 1984, och maskinläsningen av den har tappat ett par bokstäver;
  citatet återger orden som de står, med \enquote{nearly} lagat.}  Att
felet inte är ett tidsbundet Pascalfenomen visar
\textcite{Sorva2012Dissertation}, vars inventering över dokumenterade
missuppfattningar i CS1 --- den första programmeringskursen, som den
kallas i litteraturen --- bär det vidare som post 30:
\textcquote[360]{Sorva2012Dissertation}{Adjacent code executes within loop}.
Det är just den aspekten \cref{sec:indrag} varierar, och i Python är svaret
indraget.

\paragraph{Vad som måste stå efter \texttt{in}, och att räknaren måste
sättas.}  \textcite{GuoMarkelZhang2020} samlade nästan 17\,000 förklaringar
som nybörjare själva skrivit på webbplatsen Python Tutor, och katalogiserade
dem.  Bland dem finns båda felen i en och samma mening:
\textcquote[4]{GuoMarkelZhang2020}{Some learners wrote code like
\mkbibquote{while i <= 100} or \mkbibquote{for i in 100}, which read in
English like i will iterate up to 100. The former fails because i is never
initialized, and the latter fails because for-loops need a collection to
iterate over}.  Felet att skriva ett ensamt tal efter
\mintinline{python}{in} är precis det \cref{ex:hundra} ställer som fråga,
med Pythons eget felmeddelande som svar.  Felet att aldrig sätta räknaren
tas upp i \cref{sec:while}, men
inte som en kontrast: där varieras bara om kroppen \emph{ändrar} räknaren,
och att den också måste \emph{sättas} innan slingan börjar sägs i meningen
som avslutar avsnittet.  Att skriva ut även det som en kontrast vore möjligt,
och är noterat som en förbättring.

\paragraph{Villkoret som aldrig ändras.}  Belagt i studenternas egen kod.
\textcite{Bosse2021Antipatterns} katalogiserade 41 antimönster ur 95
missuppfattningar i riktiga CS1-inlämningar, och ett av dem är slingan
vars styrvariabel aldrig ändras: \textcquote[111]{Bosse2021Antipatterns}{The
program will go into an infinite loop because, since the control variable is
never changed, the result of the condition that controls the repetition will
not change either}.  Det är felet i \cref{ex:while-utan}, formulerat av dem
som sett det i inlämningarna.

\paragraph{Upprepning där en förgrening räcker.}  Samma katalog, ett annat
antimönster: \textcquote[79]{Bosse2021Antipatterns}{Using repetition where
only selection would be required will cause wrong results}.  Lägg märke till
formuleringen --- felet är att \emph{resultatet blir fel}, inte att
programmet hänger sig.  Det är därför talen i
\cref{ex:rabatt} är valda så att båda programmen terminerar.

\paragraph{Att slingor bär missuppfattningar alls.}  Att slingor, och
särskilt nästlade slingor, är en dokumenterad svårighet, och att
missuppfattningarna hänger ihop med hur slingan är \emph{representerad},
rapporterar \textcite{Eckert2022Loops}:
\textcquote{Eckert2022Loops}{For novice programmers loops generally and
nested loops specifically pose big difficulties. This paper identifies
common student misconceptions about loops in a first-year C++ introductory
computer science course and the relation of these misconceptions to
representations}.  Det motiverar både
\cref{sec:nastlade} och generaliseringen i \cref{sec:generalisering}, där
ytformen varieras och innebörden hålls fast.

\section{Diskussion och begränsningar}
\label{sec:miss-diskussion}

Omläsningen ändrade materialet på två ställen, och det är den intressanta
delen av den här bilagan.

\paragraph{Ett påstående fick skrivas om.}  Kursens forskningsunderlag
beskriver felet i \cref{sec:while} som att studenten tror att
\emph{slingvariabeln räknas upp av sig själv}.  Guo med flera säger något
smalare: att variabeln \enquote{is never initialized}.  Om automatisk
uppräkning står ingenting.  Materialets formulering är därför delad i två:
att räknaren måste \emph{sättas} före slingan beläggs med Guo med flera, och
att den måste \emph{ändras} i kroppen med Bosse med fleras antimönster ---
RS2 i deras numrering --- som säger just det.  Ingen av källorna får bära
hela påståendet ensam.

\paragraph{Ett påstående saknar källa och sägs inte.}  Missuppfattningen i
\cref{sec:for} --- att det som upprepas är den utskrivna texten och inte
kroppen --- går inte att belägga ordagrant.  Formuleringen finns inte i
\textcite{Sleeman1984}; det närmaste är felet \enquote{Data-driven Looping}
på s.~16, där \textcquote[16]{Sleeman1984}{Several students indicated that
the number of numbers in the data determined the times a loop was executed},
vilket handlar om hur \emph{många} varv slingan gör, inte om vad som
upprepas.  Att koppla ihop de två vore vår tolkning, inte källans.
Materialet tillskriver därför inte Sleeman med flera den formuleringen, och
säger uttryckligen att den saknar ordagrant stöd.

\paragraph{Övriga begränsningar.}  Tre stycken.  Eckert med fleras artikel
är bara läst i sammanfattning: IEEE\ Xplore ger inte fulltexten och någon
öppen kopia hittades inte, så den citeras bara för det sammanfattningen
själv säger.  Sleemans studie är från 1984 och gäller Pascal; att felet
återkommer i en inventering från 2012 talar för att det inte är språkbundet,
men den inventeringen är en sekundärkälla.  Och två källor som
forskningsunderlaget använder --- Sekiya och Yamaguchi om förgreningar inuti
slingor, och Kumar Veerasamy med flera om att följa kod --- gick inte att nå
i fulltext den här dagen (ACM respektive universitetsarkivet svarade med
robotspärr), så de citeras inte här.  Det är inte samma sak som att de
skulle vara fel; det betyder bara att vi inte har läst dem.

\section{Slutsats}
\label{sec:miss-slutsats}

Ja, fem av de sex påståendena är belagda i primärkällor, och det sjätte
sägs inte längre.  Kod intill slingan, kravet på något iterbart efter
\mintinline{python}{in}, räknaren som aldrig sätts, styrvariabeln som aldrig
ändras och upprepningen som borde varit en förgrening är alla dokumenterade
i studier av vad nybörjare faktiskt skriver.  Att det upprepade skulle vara
den utskrivna texten är däremot inget vi kan belägga i den formen, och
materialet påstår det inte; kontrasten i \cref{sec:for} står ändå kvar,
eftersom den gör ombindningen av slingvariabeln synlig och det behöver alla
se.  Lärdomen för den som själv ska kontrollera ett påstående är den i
\cref{sec:miss-diskussion}: en ärvd fotnot är inte ett belägg förrän du
själv har läst stället den pekar på --- och när du har läst det får
påståendet formuleras precis så starkt som stället tillåter, inte starkare.

\section{Fortsatt arbete}
\label{sec:miss-fortsatt}

\Cref{sec:miss-diskussion} räknar upp vad som begränsar svaret; det här
avsnittet säger vad som ska göras åt det.  Tre texter ska läsas i
fulltext.  Eckert med fleras artikel är hittills bara läst i
sammanfattning\autocite{Eckert2022Loops}, och just den skulle säga hur
vanliga missuppfattningarna om slingor var och i vilken grupp --- den
är den enda av kapitlets källor som räknar dem.  De två arbeten som
kursens forskningsunderlag använder men som inte gick att nå ska hämtas
genom ett bibliotek: Sekiya och Yamaguchi om förgreningar inuti slingor,
och Kumar Veerasamy med flera om att följa kod.  Det senare bär
underlagets påstående om att följa ett program steg för steg, som
materialet lutar sig mot utan att det är kontrollerat här.

För ett säkrare svar saknas framför allt siffror på nybörjare i Python.
Felen är \emph{dokumenterade} --- någon har sett dem --- men bara ett av
dem har en frekvens knuten till sig, och den kommer ur intervjuer med
gymnasieelever som skrev Pascal 1984.  Det som skulle göra svaret
säkrare är en räkning av hur ofta vart och ett av felen förekommer i
pythonnybörjares egen kod, gärna i den här kursens egna inlämningar, och
en studie där nybörjare får tänka högt medan de läser en slinga, som
prövar det sjätte påståendet direkt: tror någon nybörjare verkligen att
det som upprepas är den utskrivna texten?  Utfallet skulle synas i
materialet: höll formuleringen kunde \cref{sec:for} säga den rent ut, och
visade en räkning att något av felen är sällsynt hos pythonnybörjare
kunde den kontrast som är byggd för det lämna plats åt ett vanligare.
